The audited 50-cell rule fails joint interval evaluation despite reassuring
point-estimate recovery. A larger simulation budget distinguishes three
coverage failures from one passing condition, while excess null rejection
and insufficient detection occur in all four. Welch intervals improve coverage
but do not repair the reporting boundary. The non-monotone detection curve
also becomes interpretable only beside the changing false-positive rate.

A reporting threshold must be assessed with its actual interval, explicit
target and null and alternative operating characteristics. The synthetic
external-control example adds that the target must name the population over
which effects are averaged. These results establish failures and delimit one
attempted repair; they do not validate a replacement rule or a patient world
model.
